\section{Packed EO-LiteBullet Proof}
\label{sec:packedeolitebullet}

Our inner-product proof builds on LiteBullet
~\cite[Algorithm~1]{asiacrypt/CongCYY25}, which is recalled in
Appendix~\ref{appx:litebullet}. LiteBullet recursively folds the left and
right halves of a vector. Such folds access distant coordinates and, on our
bucketed layout, can combine stored entries with conceptual padding. We use a
low-to-high even--odd (EO) schedule instead and fuse its first $\tau$ rounds
into one packed reduction. This reduction consumes the low-bit-width input in
contiguous blocks and produces full-field vectors only after reducing their
length by a factor of $\kappa=2^\tau$.

Consider the claim
\[
    c=\langle\bm a,\bm b\rangle,
    \qquad \bm a,\bm b\in\Field^K,
\]
where $K=2^\mu$, $\bm a$ is stored at low bit width, and $\bm b$ is a public
full-field vector. The entries of $\bm a$ are interpreted canonically as field
elements only when they are consumed. Section~\ref{sec:bucketpackedeolitebullet}
shows how the BVA and TBV representations execute the same reduction without
materializing either padded vector.

\subsection{Packed Reduction}

Let $\kappa=2^\tau\le K$. View $\bm a$ and $\bm b$ as the logical matrices
\[
    \bm A
    :=[\bm a_{(1)}|\cdots|\bm a_{(\kappa)}],
    \qquad
    \bm B
    :=[\bm b_{(1)}|\cdots|\bm b_{(\kappa)}],
    \qquad
    \bm A,\bm B\in\Field^{(K/\kappa)\times\kappa},
\]
whose entries are
\begin{equation}
\label{eq:packed-matrix-layout}
    A_{i,s}:=a_{(i-1)\kappa+s},
    \qquad
    B_{i,s}:=b_{(i-1)\kappa+s},
    \qquad i\in[K/\kappa],\ s\in[\kappa].
\end{equation}
Each matrix row therefore contains one contiguous $\kappa$-element block of
the original vector. The matrices are logical views; the prover does not
materialize their columns as separate arrays.

Define the logical cross-product matrix
\[
    \bm D:=\bm A^\top\bm B\in\Field^{\kappa\times\kappa},
    \qquad
    d_{ss'}:=D_{s,s'}
    =\langle\bm a_{(s)},\bm b_{(s')}\rangle.
\]
Its trace is the claimed inner product:
\begin{equation}
\label{eq:packed-trace}
    \operatorname{tr}(\bm D)
    =\sum_{s=1}^{\kappa}d_{ss}
    =\langle\bm a,\bm b\rangle
    =c.
\end{equation}

The two Laurent vector polynomials are
\begin{align*}
    \bm g_a(X)
    &:=\bm A(1,X^{-1},\ldots,X^{-(\kappa-1)})^\top
      =\sum_{s=1}^{\kappa}X^{-(s-1)}\bm a_{(s)}, \\
    \bm g_b(X)
    &:=\bm B(1,X,\ldots,X^{\kappa-1})^\top
      =\sum_{s=1}^{\kappa}X^{s-1}\bm b_{(s)}.
\end{align*}
Their inner product satisfies
\begin{align}
\label{eq:packed-laurent}
    f(X)
    &:=\langle\bm g_a(X),\bm g_b(X)\rangle \notag\\
    &=(1,X^{-1},\ldots,X^{-(\kappa-1)})
      \bm D
      (1,X,\ldots,X^{\kappa-1})^\top \notag\\
    &=\sum_{s=1}^{\kappa}\sum_{s'=1}^{\kappa}
      d_{ss'}X^{s'-s}.
\end{align}
For $h\in\{-(\kappa-1),\ldots,\kappa-1\}$, define the sum along the
$h$-th diagonal of $\bm D$ as
\[
    c_h:=\sum_{s'-s=h}d_{ss'}.
\]
Equation~\eqref{eq:packed-trace} gives $c_0=c$. Define
\[
    \bm c_+:=(c_1,\ldots,c_{\kappa-1}),
    \qquad
    \bm c_-:=(c_{-1},\ldots,c_{-(\kappa-1)}).
\]
Then
\begin{equation}
\label{eq:packed-coefficients}
    f(X)=c+\sum_{h=1}^{\kappa-1}
    \left(c_hX^h+c_{-h}X^{-h}\right).
\end{equation}

The prover sends $\bm c_+$ and $\bm c_-$, after which the verifier samples
$y_0\ChooseRandom\Field^*$. The prover computes
\begin{align*}
    \overline{\bm a}
    &:=\bm A(1,y_0^{-1},\ldots,y_0^{-(\kappa-1)})^\top
      =\sum_{s=1}^{\kappa}y_0^{-(s-1)}\bm a_{(s)}, \\
    \overline{\bm b}
    &:=\bm B(1,y_0,\ldots,y_0^{\kappa-1})^\top
      =\sum_{s=1}^{\kappa}y_0^{s-1}\bm b_{(s)}.
\end{align*}
Both vectors have length $K/\kappa$. Equation~\eqref{eq:packed-coefficients}
reduces the original claim to
\begin{align}
\label{eq:packed-eo-claim}
    \overline c
    &:=c
      +\langle\bm c_+,(y_0,y_0^2,\ldots,y_0^{\kappa-1})\rangle \notag\\
    &\qquad
      +\langle\bm c_-,(y_0^{-1},y_0^{-2},\ldots,y_0^{-(\kappa-1)})\rangle
      =\langle\overline{\bm a},\overline{\bm b}\rangle.
\end{align}
Writing $s-1$ in binary shows that this packed reduction is exactly $\tau$
consecutive EO folds with challenges
\[
    (y_0,y_0^2,\ldots,y_0^{2^{\tau-1}}).
\]
In particular, $\kappa=2$ recovers one ordinary EO fold.

The matrix representation does not require materializing $\bm D$. The prover
accumulates its $2\kappa-1$ diagonal sums while scanning the rows of $\bm A$
and $\bm B$. Equivalently, for each row $i$, it multiplies
\[
    \left(\sum_{s=1}^{\kappa}A_{i,s}X^{-(s-1)}\right)
    \left(\sum_{s=1}^{\kappa}B_{i,s}X^{s-1}\right)
\]
and adds the coefficients to $(c_{-(\kappa-1)},\ldots,c_{\kappa-1})$.
After receiving $y_0$, a second row-wise scan writes one coordinate of
$\overline{\bm a}$ and $\overline{\bm b}$ per contiguous block. If
$M(\kappa)$ denotes the cost of multiplying two degree-$(\kappa-1)$
polynomials, the packed round takes $O((K/\kappa)M(\kappa))$ field operations
and $O(K/\kappa+\kappa)$ full-field working space. It never materializes a
$K$-element full-field representation of the low-bit-width input.

\subsection{Binary EO Completion}

The packed reduction leaves vectors of length $K/\kappa$. The remaining
$\log_2(K/\kappa)$ rounds use binary EO folds, with verifier challenges
$y_1,\ldots,y_{\log_2(K/\kappa)}$. For an even-length vector, write
$\bm a_{\sf O}=(a_0,a_2,\ldots)$ and
$\bm a_{\sf E}=(a_1,a_3,\ldots)$, and define
$\bm b_{\sf O}$ and $\bm b_{\sf E}$ analogously. Given a current claim
$c'=\langle\bm a,\bm b\rangle$, the prover sends
\[
    \iota:=\langle\bm a_{\sf O},\bm b_{\sf E}\rangle,
    \qquad
    \rho:=\langle\bm a_{\sf E},\bm b_{\sf O}\rangle.
\]
In round $j$, for a verifier challenge
$y_j\ChooseRandom\Field^*$, define
\[
    \bm a':=\bm a_{\sf O}+y_j^{-1}\bm a_{\sf E},
    \qquad
    \bm b':=\bm b_{\sf O}+y_j\bm b_{\sf E}.
\]
Then
\begin{equation}
\label{eq:binary-eo-fold}
    \langle\bm a',\bm b'\rangle
    =c'+y_j\iota+y_j^{-1}\rho.
\end{equation}
Together, the packed and binary challenges define the $\mu$-dimensional MLE
point
\[
    \bm y
    :=(y_0,y_0^2,\ldots,y_0^{2^{\tau-1}},
       y_1,\ldots,y_{\log_2(K/\kappa)}).
\]
Bit reversal maps this low-to-high EO schedule to LiteBullet's left--right
schedule, and applying the same permutation to both vectors preserves their
inner product. The binary completion therefore inherits LiteBullet's
completeness and soundness.

% !TEX root = main.tex
\begin{algorithm}[!t]
\caption{\packedeolitebullet with contiguous stride-$\kappa$ blocks}
\label{alg:packedeolitebullet}
\KwIn{$K=2^\mu$; $2\le\kappa=2^\tau\le K$; $c\in\Field$;
$\bm a,\bm b\in\Field^K$}
\KwOut{reduced scalars $\hat a,\hat b$; MLE point
$\bm y\in\Field^\mu$; acceptance bit $\mathsf{ok}$}

Let $\bm a_{\rm in}$ and $\bm b_{\rm in}$ denote the original input views\;
View $\bm a$ and $\bm b$ as
$\bm A=[\bm a_{(1)}|\cdots|\bm a_{(\kappa)}]$ and
$\bm B=[\bm b_{(1)}|\cdots|\bm b_{(\kappa)}]$, where\
\Indp
$A_{i,s}=a_{(i-1)\kappa+s}$ and
$B_{i,s}=b_{(i-1)\kappa+s}$ for
$i\in[K/\kappa]$ and $s\in[\kappa]$\;
\Indm
Let $\bm D:=\bm A^\top\bm B$ denote the logical matrix with
$d_{ss'}=\langle\bm a_{(s)},\bm b_{(s')}\rangle$\;
\tcp{Accumulate the diagonal sums without materializing $\bm D$}
\For{$h\leftarrow-(\kappa-1)$ \KwTo $\kappa-1$}{
  $c_h\leftarrow\sum_{s'-s=h}d_{ss'}$\;
}
$\bm c_+\leftarrow(c_1,\ldots,c_{\kappa-1})$; \quad
$\bm c_-\leftarrow(c_{-1},\ldots,c_{-(\kappa-1)})$\;
$\PtoV$: $(\bm c_+,\bm c_-)$; \quad
$\VtoP$: $y_0\ChooseRandom\Field^*$\;
$\overline{\bm a}\leftarrow
\bm A(1,y_0^{-1},\ldots,y_0^{-(\kappa-1)})^\top$\;
$\overline{\bm b}\leftarrow
\bm B(1,y_0,\ldots,y_0^{\kappa-1})^\top$\;
$\overline c\leftarrow
c+\sum_{h=1}^{\kappa-1}(c_hy_0^h+c_{-h}y_0^{-h})$; \quad
$\widehat c\leftarrow\overline c$\;

\For{$j\leftarrow1$ \KwTo $\log_2(K/\kappa)$}{
  EO-split
  $(\overline{\bm a}_{\sf O},\overline{\bm a}_{\sf E})
  \leftarrow\overline{\bm a}$ and
  $(\overline{\bm b}_{\sf O},\overline{\bm b}_{\sf E})
  \leftarrow\overline{\bm b}$\;
  $\iota_j\leftarrow
  \langle\overline{\bm a}_{\sf O},\overline{\bm b}_{\sf E}\rangle$; \quad
  $\rho_j\leftarrow
  \langle\overline{\bm a}_{\sf E},\overline{\bm b}_{\sf O}\rangle$\;
  \label{line:litebulletCrossIP}
  $\PtoV$: $(\iota_j,\rho_j)$; \quad
  $\VtoP$: $y_j\ChooseRandom\Field^*$\;
  $\widehat c\leftarrow
  \widehat c+y_j\iota_j+y_j^{-1}\rho_j$\;
  $\overline{\bm a}\leftarrow
  \overline{\bm a}_{\sf O}+y_j^{-1}\overline{\bm a}_{\sf E}$\;
  $\overline{\bm b}\leftarrow
  \overline{\bm b}_{\sf O}+y_j\overline{\bm b}_{\sf E}$\;
  \label{line:litebulletFoldB}
}
$\PtoV$: $\hat a\leftarrow\overline{\bm a}[0]$; \quad
$\hat b\leftarrow\overline{\bm b}[0]$\;
$\bm y\leftarrow(y_0,y_0^2,\ldots,y_0^{2^{\tau-1}},
y_1,\ldots,y_{\log_2(K/\kappa)})$\;
$\mathsf{ok}\leftarrow(\hat a\hat b\EqOrNot\widehat c)$\;
\tcp{Subsequent PCS/TBV checks}
$\mathsf{ok}\leftarrow\mathsf{ok}\wedge
[\hat a\EqOrNot\tensorproj{\bm a_{\rm in}}{\bm y^{-1}}]\wedge
[\hat b\EqOrNot\tensorproj{\bm b_{\rm in}}{\bm y}]$\;
\Return $(\mathsf{ok},\hat a,\hat b,\bm y)$\;
\end{algorithm}


\begin{theorem}[Packed EO-LiteBullet]
\label{thm:packedeolitebullet}
For $K=2^\mu$, $2\le\kappa=2^\tau\le K$, and
$\Field=\mathbb F_p$, Algorithm~\ref{alg:packedeolitebullet} has perfect
completeness and soundness error at most:
\[
    \frac{2\kappa-2+2\log_2(K/\kappa)}{p-1}.
\]
Its proof contains
$2\kappa-2+2\log_2(K/\kappa)+2$ field elements.
\end{theorem}
The proof is given in Appendix~\ref{appx:proof-packed-eo-litebullet}.


\subsection{Efficient Prover with Structured Vector}
\label{sec:litebulletAccelleration}
Suppose that the public vector is tensor structured:
\[
    \bm b=\tensorspan{\bm z},
    \qquad
    \bm z=(z_1,\ldots,z_\mu)\in\Field^\mu.
\]
Let $N\le K$ denote the number of materialized entries of $\bm a$; the
remaining coordinates, if any, are conceptual zeros.  We use the convention
$\tensorspan{\emptyset}:=(1)$ and define
\[
    \bm\zeta:=\tensorspan{(z_1,\ldots,z_\tau)}\in\Field^\kappa,
    \qquad
    \bm b':=\tensorspan{(z_{\tau+1},\ldots,z_\mu)}
       \in\Field^{K/\kappa}.
\]
Because the matrix layout in Eq.~\eqref{eq:packed-matrix-layout} groups the
low-order $\tau$ index bits within each contiguous row, its columns satisfy
\[
    \bm b_{(s)}=\zeta_s\bm b',
    \qquad
    \bm B=\bm b'\bm\zeta^\top.
\]
Consequently, the logical cross-product matrix is rank one:
\begin{equation}
\label{eq:structured-cross-product}
    \bm D
    =\bm A^\top\bm B
    =(\bm A^\top\bm b')\bm\zeta^\top
    =\bm u\bm\zeta^\top,
    \qquad
    \bm u:=\bm A^\top\bm b'.
\end{equation}
Thus, the prover first computes $\bm u$ and then obtains the packed-round
coefficients from
\[
    c_h=\sum_{s'-s=h}u_s\zeta_{s'},
    \qquad h\in\{-(\kappa-1),\ldots,\kappa-1\}.
\]
This computation materializes neither $\bm B$ nor $\bm D$.  Computing
$\bm u$ requires one multiplication between a low-bit-width entry of $\bm A$
and a full-field entry of $\bm b'$ for each of the $N$ materialized entries;
the diagonal sums require one length-$\kappa$ polynomial product.  When
$\kappa$ is chosen near the ratio of the field width to the input bit width,
these $N$ mixed-width multiplications cost $O(N/\kappa)$ full-width
multiplication equivalents under the schoolbook cost model.

The inner products in Line~\ref{line:litebulletCrossIP} and the fold in
Line~\ref{line:litebulletFoldB} can also be computed without materializing
$\overline{\bm b}$.  After the packed challenge $y_0$,
\begin{align}
\label{eq:structured-packed-fold}
    \overline{\bm b}
    &=\eta_0\tensorspan{(z_{\tau+1},\ldots,z_\mu)}, \\
    \eta_0
    &:=\sum_{s=1}^{\kappa}\zeta_sy_0^{s-1}
      =\prod_{j=1}^{\tau}(1+z_jy_0^{2^{j-1}}),
\end{align}
so $\overline{\bm b}$ remains tensor structured.

At the beginning of binary round $j$, suppose
\[
    \overline{\bm b}
    =\eta_{j-1}
      \tensorspan{(z_{\tau+j},\ldots,z_\mu)}.
\]
The two cross terms are evaluated directly as
\begin{align*}
    \iota_j
    &=\eta_{j-1}z_{\tau+j}\,
      \Eval(\overline{\bm a}_{\sf O},
            (z_{\tau+j+1},\ldots,z_\mu)), \\
    \rho_j
    &=\eta_{j-1}\,
      \Eval(\overline{\bm a}_{\sf E},
            (z_{\tau+j+1},\ldots,z_\mu)).
\end{align*}
The fold then updates only the tensor scale:
\begin{equation}
\label{eq:structured-b-fold}
    \eta_j
    \leftarrow\eta_{j-1}(1+y_j z_{\tau+j}),
    \qquad
    \overline{\bm b}
    =\eta_j
      \tensorspan{(z_{\tau+j+1},\ldots,z_\mu)}.
\end{equation}
Thus, $\overline{\bm b}$ is never materialized.  Since the length of
$\overline{\bm a}$ decreases geometrically from $O(N/\kappa)$, all binary
rounds together require $O(N/\kappa)$ full-field multiplications.


\subsection{Efficient Prover with Tensor-Bucketed Vectors}

We extend the preceding tensor-structured vector to a concatenation of aligned
tensor blocks.  This representation preserves the time and space efficiency of
the structured-vector prover while supporting heterogeneous block sizes.

\begin{definition}[Tensor-Bucketed Vector]
\label{def:tbv}
Let $T\in\mathbb{N}$ and let $k_1,\ldots,k_T\in\mathbb{N}$ satisfy:
\begin{itemize}[label={},leftmargin=0pt]
    \item \textcircled{1} each $k_i$ and the logical length $K$ are powers of
    two, and $S:=\sum_{i=1}^T k_i\le K$;
    \item \textcircled{2} $k_1\ge k_2\ge\cdots\ge k_T$.
\end{itemize}
A \emph{tensor-bucketed vector (TBV)} over $\Field$ is specified by the tuple:
\[
\big(
K,\,
T,\,
\{\,k_i,\eta_i,\bm{z}_i\,\}_{i\in[T]}
\big),
\]
where $\eta_i\in\Field$ and $\bm{z}_i\in\Field^{\log_2 k_i}$,
and expands to the vector:
\[
\bm{\zeta}
=
\big(
\eta_1\cdot\tensorspan{\bm{z}_1}
~\|~
\eta_2\cdot\tensorspan{\bm{z}_2}
~\|~
\cdots
~\|~
\eta_T\cdot\tensorspan{\bm{z}_T}
~\|~
\bm 0^{K-S}
\big)
\in\Field^K .
\]
The final $K-S$ coordinates form one implicit zero suffix; they are not a
bucket and are never materialized.
The set of all such vectors is denoted by $\TBV(\Field)$.
\end{definition}

Let $K_0:=0$ and $K_i:=\sum_{j=1}^i k_j$.  Since the bucket sizes are
non-increasing powers of two, $k_i\mid k_j$ for every $j<i$, and hence
$k_i\mid K_{i-1}$.  Thus, every bucket is aligned to its own size.  Moreover,
all unit buckets form a suffix, followed only by the implicit zero suffix.

\mysmallskip
{\bf Compact representation.}
A TBV is fully specified by the tuple:
\[
    \big(
        K,\,
        T,\,
        \{\, k_i,\, \eta_i,\, \bm{z}_i\,\}_{i \in [T]}
    \big).
\]
Since $|\bm z_i|=\log_2 k_i$, for $K\ge2$ the description occupies
\[
    \Theta\!\left(T+\sum_{i=1}^T\log_2 k_i\right)
    =O(T\log K)
\]
field words, rather than $K$ field elements for the expanded vector.


\mysmallskip
\noindent
\textbf{Packed EO-LiteBullet without materialization.}

A TBV is closed under EO splitting and folding.  Let
$T_1:=|\{i:k_i=1\}|$ and $J:=T-T_1$.  One EO round halves every bucket
$i\le J$ and pairs consecutive buckets in the unit suffix.  If $T_1$ is odd,
then $S$ is odd; because $K\ge2$ is even, $S<K$, so the implicit zero suffix is
nonempty and its first coordinate pairs with the last unit bucket.  The folded
TBV therefore has the non-increasing bucket sizes
\[
  \bigl(k_1/2,\ldots,k_J/2,
  \underbrace{1,\ldots,1}_{\lceil T_1/2\rceil}\bigr).
\]
Within a non-unit bucket, the two EO halves share the remaining tensor point;
folding removes its first coordinate and updates its scale.  A pair of unit
buckets is folded by combining their scales; an unmatched final unit retains
its scale because its partner is zero.  Its active cover length is
$\lceil S/2\rceil$, and the remaining $K/2-\lceil S/2\rceil$ coordinates are
again one implicit zero suffix.  The result is therefore another TBV of length
$K/2$, computed entirely from the compact description.  Since
the packed stride-$\kappa$ reduction is $\tau$ EO folds with correlated
challenges $(y_0,y_0^2,\ldots,y_0^{2^{\tau-1}})$, packing a TBV also produces a
TBV with active cover length $\lceil S/\kappa\rceil$, without materializing an
intermediate query vector.


\begin{property}[TBV Inner Products]
\label{prop:iptbv}
\label{prop:mixiptbv}
Let $\bm\zeta\in\TBV(\Field)$ and partition
$\bm a=(\bm a^{(1)}\|\cdots\|\bm a^{(T)}\|\bm a^{(\mathsf{pad})})
\in\Field^K$ according to the bucket sizes and the final suffix.  Then
\begin{equation}
\label{eq:mixtbvip}
    \langle\bm a,\bm\zeta\rangle
    =\sum_{i=1}^T\eta_i\,
      \Eval(\bm a^{(i)},\bm z_i),
\end{equation}
so an inner product with a TBV is computed through one evaluation per bucket.
Moreover, if
$\bm\zeta'=(K,T,\{k_i,\eta_i',\bm z_i'\}_{i\in[T]})$ shares the same
bucket sizes, then
\[
    \langle\bm\zeta,\bm\zeta'\rangle
    =\sum_{i=1}^T\eta_i\eta_i'
      \prod_{j=1}^{\log_2 k_i}(1+z_{ij}z_{ij}').
\]
Both computations avoid expanding either TBV.
\end{property}

Equation~\eqref{eq:mixtbvip} computes the cross inner products in
Line~\ref{line:litebulletCrossIP}, while TBV closure implements the fold in
Line~\ref{line:litebulletFoldB}.  Across all binary rounds, processing the TBV
description takes $O(T\log K)$ time and space.  The witness vectors shrink
geometrically from length $O(N/\kappa)$, so their total inner-product work is
$O(N/\kappa)$.  In particular, replacing one tensor-structured query with a
TBV introduces no time or auxiliary-memory term proportional to $K$.
